% Źródła: Hartshorne s. 296--303 (§33); Eves s. 284, 288--291; Greenberg 2010, s. 201, 204--207, 210, 214; Wikipedia: Playfair's axiom, John Wallis, Clavius, Clairaut, Proclus.
% Uwaga: aksjomat Arystotelesa NIE jest równoważny piątemu postulatowi (zachodzi także w geometrii hiperbolicznej).

\section{Dobre, stare} % lang-pl
\section{I classici} % lang-it
\label{sekcja_dobre_stare}

W tej sekcji zbierzemy zdania, które od starożytności do XIX wieku będą podawane w miejsce piątego postulatu, każde z nazwiskiem swojego autora. % lang-pl
Przy pozostałych aksjomatach geometrii neutralnej większość z nich okaże się równoważna piątemu postulatowi; wyjątkiem będzie aksjomat Arystotelesa. % lang-pl
Dowody równoważności zostawimy czytelnikowi -- Eves \cite[s. 284, 288, 291]{eves1_1972} i Hartshorne \cite[s. 301--303]{hartshorne_2000} (zad. 33.1--33.11) proponują je jako zadania. % lang-pl
In questa sezione raccoglieremo gli enunciati che dall'antichità all'Ottocento saranno proposti al posto del quinto postulato, ciascuno con il nome del suo autore. % lang-it
Assumendo gli altri assiomi della geometria neutrale, la maggior parte di essi risulterà equivalente al quinto postulato; farà eccezione l'assioma di Aristotele. % lang-it
Le dimostrazioni delle equivalenze le lasceremo al lettore -- Eves \cite[p. 284, 288, 291]{eves1_1972} e Hartshorne \cite[p. 301--303]{hartshorne_2000} (es. 33.1--33.11) le propongono come esercizi. % lang-it

\begin{proposition}
[aksjomat Arystotelesa] % lang-pl
[assioma di Aristotele] % lang-it
	Dla każdego kąta ostrego, każdego jego ramienia i każdego odcinka $AB$ istnieje na tym ramieniu punkt $Y$ taki, że jeśli $X$ jest spodkiem prostopadłej z $Y$ na drugie ramię, to $YX > AB$. % lang-pl
	Per ogni angolo acuto, ogni suo lato e ogni segmento $AB$ esiste su quel lato un punto $Y$ tale che, se $X$ è il piede della perpendicolare da $Y$ all'altro lato, allora $YX > AB$. % lang-it
\index{aksjomat!Arystotelesa}%
\end{proposition}
Arystoteles wypowie to zdanie w księdze I traktatu \emph{De caelo}, a Proklos użyje go w V wieku w nieudanej próbie dowodu piątego postulatu \cite[s. 204--205]{greenberg_2010}. % lang-pl
Zdanie wynika z piątego postulatu i zachodzi w każdej płaszczyźnie hiperbolicznej, a zawodzi na sferze i w geometrii eliptycznej; wynika też z aksjomatu Archimedesa, ale nie odwrotnie \cite[s. 205, 210]{greenberg_2010}. % lang-pl
Greenberg nazywa twierdzeniem Proklosa fakt, że w płaszczyźnie, w której suma kątów każdego trójkąta jest równa $180^\circ$, aksjomat Arystotelesa jest równoważny aksjomatowi równoległych \cite[s. 207]{greenberg_2010}. % lang-pl
Aristotele enuncerà questo principio nel libro I del trattato \emph{De caelo}, e Proclo lo userà nel V secolo nel suo fallito tentativo di dimostrare il quinto postulato \cite[p. 204--205]{greenberg_2010}. % lang-it
L'enunciato discende dal quinto postulato e vale in ogni piano iperbolico, mentre cade sulla sfera e nella geometria ellittica; discende anche dall'assioma di Archimede, ma non viceversa \cite[p. 205, 210]{greenberg_2010}. % lang-it
Greenberg chiama teorema di Proclo il fatto che in un piano in cui la somma degli angoli di ogni triangolo è $180^\circ$ l'assioma di Aristotele equivale all'assioma delle parallele \cite[p. 207]{greenberg_2010}. % lang-it

\begin{proposition}
[lemat Proklusa] % lang-pl
[lemma di Proclo] % lang-it
	Prosta nie może przecinać tylko jednej z dwóch prostych równoległych. % lang-pl
	Una retta non può intersecare soltanto una di due rette parallele. % lang-it
\index{lemat!Proklusa}%
\end{proposition}

Proklos sformułuje to zdanie w V wieku w komentarzu do twierdzenia I.29 i pokaże, że wynika ono z aksjomatu Arystotelesa; równoważność lematu z aksjomatem Playfaira to zadanie 33.1 u Hartshorne'a \cite[s. 297, 301]{hartshorne_2000}. % lang-pl
Rozumowanie Proklosa przytoczy bez uwag krytycznych w 1575 roku F.~Commandino w swoim wydaniu Euklidesa \cite[s. 297]{hartshorne_2000}. % lang-pl
Proclo formulerà questo enunciato nel V secolo nel commento alla proposizione I.29 e mostrerà che discende dall'assioma di Aristotele; l'equivalenza del lemma con l'assioma di Playfair è l'esercizio 33.1 di Hartshorne \cite[p. 297, 301]{hartshorne_2000}. % lang-it
Il ragionamento di Proclo sarà riportato senza commenti critici nel 1575 da F.~Commandino nella sua edizione di Euclide \cite[p. 297]{hartshorne_2000}. % lang-it

\begin{proposition}
[aksjomat Claviusa] % lang-pl
[assioma di Clavio] % lang-it
	Zbiór punktów leżących po jednej stronie danej prostej i jednakowo od niej odległych jest prostą. % lang-pl
	L'insieme dei punti che stanno da una stessa parte di una retta data e hanno da essa la stessa distanza è una retta. % lang-it
\index{aksjomat!Claviusa}%
\end{proposition}

Clavius uzna to za oczywiste w swoim wydaniu Euklidesa z 1574 roku, opierając się na starszych rozważaniach Proklosa i Posejdoniosa; w płaszczyźnie hiperbolicznej krzywe równoodległe od prostej nie są jednak prostymi \cite[s. 214]{greenberg_2010}. % lang-pl
Hartshorne \cite[s. 298--299]{hartshorne_2000} opowiada, że tego niejawnego założenia użyje też Andrea Tacquet w bardzo popularnym wydaniu \emph{Elementów} z 1654 roku, definiując proste równoległe jako wszędzie jednakowo od siebie odległe; jego zdanie o równoodległych to prawie to samo, co postulat równoległych (zad. 33.7). % lang-pl
Eves \cite[s. 284, 291]{eves1_1972} (zad. 15) podaje bliskie zdanie bez nazwiska: istnieją dwie proste leżące w jednej płaszczyźnie i wszędzie jednakowo od siebie odległe. % lang-pl
Clavio lo riterrà evidente nella sua edizione di Euclide del 1574, appoggiandosi a considerazioni più antiche di Proclo e Posidonio; nel piano iperbolico, tuttavia, le curve equidistanti da una retta non sono rette \cite[p. 214]{greenberg_2010}. % lang-it
Hartshorne \cite[p. 298--299]{hartshorne_2000} racconta che questa ipotesi implicita la userà anche Andrea Tacquet nell'assai diffusa edizione degli \emph{Elementi} del 1654, definendo parallele le rette ovunque equidistanti; il suo enunciato sulle equidistanti è quasi lo stesso del postulato delle parallele (es. 33.7). % lang-it
Eves \cite[p. 284, 291]{eves1_1972} (es. 15) riporta un enunciato affine senza nome: esistono due rette complanari ovunque equidistanti tra loro. % lang-it

\begin{proposition}
[aksjomat Clairauta] % lang-pl
[assioma di Clairaut] % lang-it
	Niech $AC$ i $BD$ będą równymi odcinkami prostopadłymi do odcinka $AB$. % lang-pl
	Wtedy kąty przy $C$ i $D$ są proste, czyli $ABCD$ jest prostokątem. % lang-pl
	Siano $AC$ e $BD$ segmenti uguali perpendicolari al segmento $AB$. % lang-it
	Allora gli angoli in $C$ e $D$ sono retti, cioè $ABCD$ è un rettangolo. % lang-it
\index{aksjomat!Clairauta}%
\end{proposition}

Alexis Claude Clairaut (1713--1765), uczestnik wyprawy do Laponii, wyda w 1741 roku \emph{Éléments de géométrie}, w których zamiast definicji i postulatów, które „nie dają czytelnikowi nic prócz suchości'', będzie prowadzić ucznia przez pomiary ziemi. % lang-pl
Prosta konstrukcja prostokąta -- w końcach odcinka $AB$ wystawić prostopadłe równej długości i połączyć ich końce -- ukrywa założenie, które Hartshorne nazywa aksjomatem Clairauta \cite[s. 299]{hartshorne_2000}; jest on równoważny aksjomatowi Claviusa (zad. 33.9, s. 302). % lang-pl
Alexis Claude Clairaut (1713--1765), partecipante alla spedizione in Lapponia, pubblicherà nel 1741 gli \emph{Éléments de géométrie}, in cui invece di definizioni e postulati, che «non offrono al lettore nulla se non aridità», guiderà l'allievo attraverso le misure del terreno. % lang-it
La semplice costruzione di un rettangolo -- innalzare agli estremi del segmento $AB$ perpendicolari di uguale lunghezza e congiungerne gli estremi -- nasconde un'ipotesi che Hartshorne chiama assioma di Clairaut \cite[p. 299]{hartshorne_2000}; esso è equivalente all'assioma di Clavio (es. 33.9, p. 302). % lang-it
\index[persons]{Clairaut, Alexis Claude}%

\begin{proposition}
[aksjomat Simsona] % lang-pl
[assioma di Simson] % lang-it
	Prosta nie może najpierw zbliżyć się do innej prostej, a potem oddalić od niej, zanim ją przetnie; podobnie nie może najpierw oddalić się od innej prostej, a potem do niej zbliżyć; nie może też przez jakiś czas zachowywać stałą odległość od innej prostej, a potem się do niej zbliżyć lub od niej oddalić. % lang-pl
	Una retta non può prima avvicinarsi a un'altra retta e poi allontanarsene prima di intersecarla; analogamente non può prima allontanarsi da un'altra retta e poi avvicinarsi ad essa; né può per un tratto mantenere la stessa distanza da un'altra retta e poi avvicinarsi o allontanarsi. % lang-it
\index{aksjomat!Simsona}%
\end{proposition}
Robert Simson (1687--1768), profesor uniwersytetu w Glasgow, wyda w 1756 roku łacińsko-angielskie wydanie \emph{Elementów}, które przejdzie około trzydziestu wydań; chciał „przywrócić Euklidesowi'' to, co wydawcy zniekształcili. % lang-pl
O piątym postulacie napisze, że nie jest on właściwie umieszczony wśród aksjomatów, bo nie jest oczywisty, ale da się go udowodnić -- i wprowadzi swój aksjomat; korzystając z niego, a niejawnie z aksjomatu Archimedesa, udowodni (poprawnie) pięć twierdzeń, z których ostatnie to piąty postulat \cite[s. 299--300]{hartshorne_2000}. % lang-pl
Aksjomat Simsona jest równoważny aksjomatowi Claviusa (zad. 33.10) \cite[s. 302]{hartshorne_2000}. % lang-pl
Robert Simson (1687--1768), professore all'Università di Glasgow, pubblicherà nel 1756 un'edizione latina e inglese degli \emph{Elementi}, che avrà circa trenta edizioni; voleva «restituire a Euclide» ciò che gli editori avevano alterato. % lang-it
Sul quinto postulato scriverà che non è collocato bene tra gli assiomi, perché non è evidente, ma si può dimostrare -- e introdurrà il suo assioma; servendosi di esso, e implicitamente dell'assioma di Archimede, dimostrerà (correttamente) cinque proposizioni, l'ultima delle quali è il quinto postulato \cite[p. 299--300]{hartshorne_2000}. % lang-it
L'assioma di Simson è equivalente all'assioma di Clavio (es. 33.10) \cite[p. 302]{hartshorne_2000}. % lang-it
\index[persons]{Simson, Robert}%

\begin{proposition}
[aksjomat Playfaira] % lang-pl
[assioma di Playfair] % lang-it
	Przez punkt leżący poza prostą przechodzi co najwyżej jedna prosta do niej równoległa (równoważnie: dwie proste przecinające się nie mogą być obie równoległe do tej samej prostej). % lang-pl
	Per un punto esterno a una retta passa al più una retta a essa parallela (equivalentemente: due rette incidenti non possono essere entrambe parallele a una stessa retta). % lang-it
\index{aksjomat!Playfaira}%
\end{proposition}

Treść aksjomatu Playfaira poda szkocki matematyk John Playfair (1748--1819), profesor filozofii przyrody w Edynburgu, w podręczniku \emph{Elements of Geometry} z 1795 roku, będącym nowym wydaniem pierwszych sześciu ksiąg Simsona \cite[s. 300]{hartshorne_2000}. % lang-pl
L'enunciato dell'assioma di Playfair fu formulato dal matematico scozzese John Playfair (1748--1819), professore di filosofia naturale a Edimburgo, nel manuale \emph{Elements of Geometry} del 1795, nuova edizione dei primi sei libri di Simson \cite[p. 300]{hartshorne_2000}. % lang-it
% % https://en.wikipedia.org/wiki/Playfair%27s_axiom
\index[persons]{Playfair, John}%
Greenberg \cite[s. 201]{greenberg_2010} zauważa, że zwykle podaje się je błędnie, dodając żądanie istnienia równoległej -- a Hilbert przyjmie właśnie wersję ``nie więcej niż jedna''. % lang-pl
Zdanie to ma jednak starszą historię: Proklos opisze je w komentarzu do twierdzenia I.31, a William Ludlam zapisze w 1785 roku, że dwie proste przecinające się nie mogą być obie równoległe do trzeciej -- i Playfair przyzna, że korzysta z uproszczeń Ludlama i innych. % lang-pl
Hilbert wybierze właśnie ten aksjomat do swoich \emph{Grundlagen der Geometrie} (1899). % lang-pl
Twierdzenia I.27 i I.28 gwarantują istnienie co najmniej jednej równoległej \cite[s. 284]{eves1_1972}, więc aksjomat Playfaira mówi, że jest \emph{dokładnie} jedna. % lang-pl
Greenberg \cite[p. 201]{greenberg_2010} osserva che di solito lo si cita erroneamente aggiungendo l'esistenza della parallela -- e Hilbert adotterà proprio la versione «non più di una». % lang-it
L'enunciato ha però una storia più antica: Proclo lo descriverà nel commento alla proposizione I.31, e William Ludlam scriverà nel 1785 che due rette incidenti non possono essere entrambe parallele a una terza -- e Playfair riconoscerà di servirsi delle semplificazioni di Ludlam e di altri. % lang-it
Hilbert sceglierà proprio questo assioma per i suoi \emph{Grundlagen der Geometrie} (1899). % lang-it
Le proposizioni I.27 e I.28 garantiscono l'esistenza di almeno una parallela \cite[p. 284]{eves1_1972}, dunque l'assioma di Playfair dice che ce n'è \emph{esattamente} una. % lang-it
\index[persons]{Ludlam, William}%

\begin{proposition}
[aksjomat Wallisa] % lang-pl
[assioma di Wallis] % lang-it
	Dla każdego trójkąta $ABC$ i odcinka $DE$ istnieje trójkąt podobny $A'B'C'$ (o tych samych kątach) z bokiem $A'B' \ge DE$. % lang-pl
	Per ogni triangolo $ABC$ e ogni segmento $DE$ esiste un triangolo simile $A'B'C'$ (con gli stessi angoli) con lato $A'B' \ge DE$. % lang-it
\index{aksjomat!Wallisa}%
\end{proposition}

John Wallis (1616--1703), profesor geometrii w Oksfordzie od 1649 roku, sformułuje to zdanie w 1663 roku po lekturze Sadr ad-Dina. % lang-pl
Zasada, że do każdej figury istnieje figura podobna dowolnej wielkości, prowadzi do dowodu piątego postulatu (zad. 33.5); Hartshorne \cite[s. 301--302]{hartshorne_2000} pokazuje, że w nieArchimedesowej geometrii nieeuklidesowej istnieją trójkąty podobne o różnych rozmiarach, ale aksjomat Wallisa w niej zawodzi. % lang-pl
Eves \cite[s. 284, 291 (zad. 14)]{eves1_1972} podaje wersję: istnieją dwa trójkąty podobne i nieprzystające. % lang-pl
John Wallis (1616--1703), professore di geometria a Oxford dal 1649, formulerà questo enunciato nel 1663 dopo aver letto Sadr al-Din. % lang-it
Il principio che per ogni figura esiste una figura simile di grandezza arbitraria porta a una dimostrazione del quinto postulato (es. 33.5); Hartshorne \cite[p. 301--302]{hartshorne_2000} mostra che in una geometria non euclidea non archimedea esistono triangoli simili di dimensioni diverse, ma l'assioma di Wallis in essa cade. % lang-it
Eves \cite[p. 284, 291 (es. 14)]{eves1_1972} riporta la versione: esistono due triangoli simili e non congruenti. % lang-it

\begin{proposition}
[aksjomat Bolyi] % lang-pl
[assioma di Bolyai] % lang-it
	Przez trzy punkty niewspółliniowe przechodzi okrąg. % lang-pl
	Per tre punti non allineati passa una circonferenza. % lang-it
\index{aksjomat!Bolyi}%
\end{proposition}
Zdanie to zaproponuje Farkas Bolyai, ojciec Jánosa; wynika ono z aksjomatu Playfaira, a także go implikuje (zad. 33.11) \cite[s. 302--303]{hartshorne_2000}. % lang-pl
Questo enunciato lo proporrà Farkas Bolyai, padre di János; discende dall'assioma di Playfair e lo implica (es. 33.11) \cite[p. 302--303]{hartshorne_2000}. % lang-it

\begin{proposition}
[aksjomat Legendre'a] % lang-pl
[assioma di Legendre] % lang-it
	Przez każdy punkt wewnątrz kąta mniejszego niż $60^\circ$ można poprowadzić prostą przecinającą oba ramiona kąta. % lang-pl
	Per ogni punto interno a un angolo minore di $60^\circ$ si può tracciare una retta che interseca entrambi i lati dell'angolo. % lang-it
\index{aksjomat!Legendre'a}%
\end{proposition}

Ta wersja pochodzi z rozumowania Legendre'a, o którym opowiemy w jego własnej sekcji; Eves \cite[s. 284, 290 (zad. 12)]{eves1_1972}. % lang-pl
Questa versione nasce da un ragionamento di Legendre, di cui parleremo nella sua sezione; Eves \cite[p. 284, 290 (es. 12)]{eves1_1972}. % lang-it

Na liście Evesa \cite[s. 284]{eves1_1972} są jeszcze dwa zdania bez nazwiska: istnieje trójkąt o sumie kątów równej dwóm kątom prostym oraz pole trójkąta nie ma ograniczenia z góry. % lang-pl
Ostatnie z nich okaże się ważne w geometrii Łobaczewskiego. % lang-pl
Nell'elenco di Eves \cite[p. 284]{eves1_1972} ci sono ancora due enunciati senza nome: esiste un triangolo con somma degli angoli uguale a due angoli retti e l'area di un triangolo non ha limite superiore. % lang-it
Quest'ultimo si rivelerà importante nella geometria di Lobachevskij. % lang-it

% TODO: https://www.deltami.edu.pl/2016/09/inne-swiaty-inne-geometrie/
